% page 152 du fac-simile (image p-151 du scan)
% bloc 152.4 x 233.3 mm ; marge gauche 8.0 mm
\pgc{-12.700mm}{0.000mm}{
\l{\cel{3}142}
\l{}
\l{}
\l{\sou{erizipelo}. (\sou{patol}.)Influro surfacala di la pelo, kun ten-}
\l{\cel{3}seso. - EFIS.}
\l{}
\l{\sou{ermino}. (\sou{zool.}) Martro blanka, mamifero digitigrada di qua}
\l{\cel{3}la pelo, kovrita subventre da pili blanka tre tenua,}
\l{\cel{3}esas furo tre chera. - L.\cel{1}\sou{putorius ermineus}. - DEFIS.}
\l{}
\l{\sou{ermito}. (\sou{religio}). Solitaro, retretinta a loko dezerta,}
\l{\cel{3}por ibe exercar su a pieso. - DEFIRS.}
\l{}
\l{\sou{erodar}. (\sou{trans.}) Rodar per efiko korodiva. - DEFIS.}
\l{}
\l{\sou{erorar}. (\sou{netrans}.)I. Forirar de lo exakta. - II. (\sou{pri dok}-}
\l{\cel{3}\sou{trino, opiniono}).Ne esar konforma a lo exakta, a lo}
\l{\cel{3}vera. - EFIS.}
\l{}
\l{\sou{erotika}. Qua koncernas amoro. - DEFIRS.}
\l{}
\l{\sou{erste}. Ne ante (ta o ca dato, ta o ca evento). - D.}
\l{}
\l{\sou{erudar}. (\sou{trans., pri}).Exercar pri la cienco di la dokumen-}
\l{\cel{3}ti qui koncernas ta o ca ek la savaji homala. - EFIS.}
\l{}
\l{\sou{eruptar}. (\sou{netrans}.)(\sou{pri kozo}). Ekirar, desintrikar su bruske}
\l{\cel{3}de to quo kontenas lu. - DEFIS.}
\l{}
\l{\sou{ervo}. (\sou{bot.}) Genero de leguminosi, di qua la speco tipa esas}
\l{\cel{3}la lentilio. - L.\cel{1}\sou{ervum}. - DEFISL.}
\l{}
\l{\sou{-es-}.\cel{2}Sufixo qua indikas la qualeso, la estado, la stando.}
\l{}
\l{\sou{esar}.\cel{2}(\sou{netrans.}) Verbo qua ligas la atributo a la subjekto}
\l{\cel{3}di la propoziciono, afirmante la qualeso di olca. -EFIRS.}
\l{}
\l{\sou{esamo}. I. Grupo de abeli, de vespi, qui vivas socie. - II.}
\l{\cel{3}Grupo multa-membro de insekti. -(\sou{Anke metafore}). FL.}
\l{}
\l{\sou{esayo}. Libro en qua la autoro traktas surfacale temo, sen}
\l{\cel{3}vizir a traktar lu profunde. - DEFS.}
\l{}
\l{\sou{esbosar}. (\sou{trans.}) Komencar laboro, verko, indikante lua}
\l{\cel{3}formi, kolori, e c., sen par-facar ula ek li. - FI.}
\l{}
\l{\sou{esenco}. I. (\sou{filoz. e teol}.)Quo konstitucas la fundo di la}
\l{\cel{3}ento. - II. Quo konstitucas la naturo propra di kozo.}
\l{\cel{3}- III. Quo konstitucas la parto maxim importoza di kozo.}
\l{\cel{3}- IV. Quo kontenas la efiko partikulara ye substanco.}
\l{\cel{3}- DEFIRS.}
\l{}
\l{\sou{esforcar}. (\sou{netrans}.)Uzar sua vigoro por rezistar o por vin-}
\l{\cel{3}kar rezistanto. - FIS.}
}
